    Then we have
    \[
    \prod_{j = 1}^{k} I_{1, j, m_{j}} \subseteq \bigcup_{\alpha \in A_{\bm{m}}} G_{\alpha}
    \quad \text{for all } \bm{m} \in S_{k},
    \]
    where $\displaystyle \bigcup_{\alpha \in A_{\bm{m}}} G_{\alpha}$ is open by \Cref{thm:ArbitraryUnionOfOpenSetsIsOpen}.
    Then we have
    \[
    I_{1} \subseteq \bigcup_{\bm{m} \in S_{k}} \paren*{\bigcup_{\alpha \in A_{\bm{m}}} G_{\alpha}}
    = \bigcup_{\alpha \in A} G_{\alpha},
    \]
    where $A \coloneqq \displaystyle \bigcup_{\bm{m} \in S_{k}} A_{\bm{m}}$ is finite by
    \Cref{thm:FiniteUnionOfFiniteSetsIsFinite}. Thus $I_{1}$ has a finite subcover, which is a contradiction.
    Therefore, there exists $\bm{m_{1}} \in S_{k}$ such that
    \[
    \prod_{j = 1}^{k} I_{1, j, \paren*{m_{1}}_{j}} \text{ has no finite subcover}.
    \]
    Let $I_{2}$ be a $k$-cell in $\R^{k}$ defined by
    \[
    I_{2} \coloneqq \prod_{j = 1}^{k} \bracket{a_{2, j}, b_{2, j}} \subseteq \R^{k},
    \]
    where $a_{2, j}, b_{2, j} \in \R$ is defined by
    \begin{itemize}
        \item $a_{2, j} \coloneqq \begin{dcases}
            a_{1, j} & \text{if } \paren*{m_{1}}_{j} = 1 \\
            \frac{a_{1, j} + b_{1, j}}{2} & \text{if } \paren*{m_{1}}_{j} = 2
        \end{dcases}$
        \item $b_{2, j} \coloneqq \begin{dcases}
            \frac{a_{1, j} + b_{1, j}}{2} & \text{if } \paren*{m_{1}}_{j} = 1 \\
            b_{1, j} & \text{if } \paren*{m_{1}}_{j} = 2
        \end{dcases}$
    \end{itemize}
    for all $j = 1, 2, \dots, k$. Then, for each $j = 1, 2, \dots, k$, we have
    \[
    \bracket{a_{2, j}, b_{2, j}} = I_{1, j, \paren*{m_{1}}_{j}} \subseteq \bracket{a_{1, j}, b_{1, j}},
    \]
    which implies that
    \[
    I_{2} = \prod_{j = 1}^{k} I_{1, j, \paren*{m_{1}}_{j}} \subseteq I_{1}.
    \]
    Let $I_{2, j, m}$ be an interval in $\R$ defined by
    \[
    I_{2, j, m} \coloneqq \begin{dcases}
        \bracket*{a_{2, j}, \frac{a_{2, j} + b_{2, j}}{2}} & \text{if } m = 1 \\
        \bracket*{\frac{a_{2, j} + b_{2, j}}{2}, b_{2, j}} & \text{if } m = 2
    \end{dcases} \quad \text{for all } j = 1, 2, \dots, k.
    \]
    Then we have
    \[
    I_{2} = \bigcup_{\bm{m} \in S_{k}} \paren*{\prod_{j = 1}^{k} I_{2, j, m_{j}}}.
    \]
    We have shown that $\set*{G_{\alpha}}_{\alpha \in A}$ is an open cover of $I_{2}$, and $I_{2}$ has no finite subcover.
    Then, for each $\bm{m} \in S_{k}$, we have
    \[
    \prod_{j = 1}^{k} I_{2, j, m_{j}} \subseteq I_{2} \subseteq \bigcup_{\alpha \in A} G_{\alpha},
    \]
    which implies that $\set*{G_{\alpha}}_{\alpha \in A}$ is an open cover of
    $\displaystyle \prod_{j = 1}^{k} I_{2, j, m_{j}}$. Therefore, there exists $\bm{m_{2}} \in S_{k}$ such that
    \[
    \prod_{j = 1}^{k} I_{2, j, \paren*{m_{2}}_{j}} \text{ has no finite subcover}.
    \]
    By repeating this process, we can construct the following sequences:
    \begin{itemize}
        \item The sequence $\paren*{\bm{m}_{n}}_{n = 1}^{\infty}$ in $S_{k}$.
        \item The sequences $\paren*{a_{n, j}}_{n = 1}^{\infty}$ and
        $\paren*{b_{n, j}}_{n = 1}^{\infty}$ in $\R$ for all $j = 1, 2, \dots, k$.
        \item The sequence $\paren*{I_{n, j, m}}_{n = 1}^{\infty}$ of intervals in $\R$
        for all $j = 1, 2, \dots, k$ and $m = 1, 2$.
        \item The sequence $\paren*{I_{n}}_{n = 1}^{\infty}$ of $k$-cells in $\R^{k}$.
    \end{itemize}
    We have the following properties:
    \begin{itemize}
        \item $I_{n + 1} \subseteq I_{n}$ for all $n \in \N$.
        \item $I_{n}$ has no finite subcover for all $n \in \N$.
        \item $a_{n, j} \leq a_{n + 1, j} \leq b_{n + 1, j} \leq b_{n, j}$ for all $j = 1, 2, \dots, k$ and $n \in \N$.
        \item $b_{n + 1, j} - a_{n + 1, j} = \dfrac{1}{2}(b_{n, j} - a_{n, j}) \implies b_{n, j} - a_{n, j} =
        \dfrac{1}{2^{n - 1}}(b_{1, j} - a_{1, j})$ for all $j = 1, 2, \dots, k$ and $n \in \N$.
    \end{itemize}
    Let $\delta \coloneqq \displaystyle \sqrt{\sum_{j = 1}^{k} (b_{1, j} - a_{1, j})^{2}}$. Then we have
    \[
    \norm*{\bm{x} - \bm{y}} \leq \sqrt{\sum_{j = 1}^{k} (b_{n, j} - a_{n, j})^{2}} =
    \frac{\delta}{2^{n - 1}} \quad \text{for all } \bm{x}, \bm{y} \in I_{n} \text{ and } n \in \N.
    \]
    By \Cref{thm:NestedKCells}, we have $\displaystyle \bigcap_{n = 1}^{\infty} I_{n} \neq \vnot$.
    Let $\bm{\beta} \in \displaystyle \bigcap_{n = 1}^{\infty} I_{n}$ be given. Then there exists $\alpha_{0} \in A$
    such that $\bm{\beta} \in G_{\alpha_{0}}$. Since $G_{\alpha_{0}}$ is open, there exists $\veps \in \R^{+}$ such that
    \[
    B_{\veps}(\bm{\beta}) \subseteq G_{\alpha_{0}}.
    \]
    By \Cref{cor:ArchimedeanPropertyCorollary}, there exists $N \in \N$
    such that $\dfrac{1}{2^{N}} < \dfrac{\veps}{2\delta}$. Since $\bm{\beta} \in I_{N}$, we have
    \[
    \norm*{\bm{\beta} - \bm{\gamma}} \leq \frac{\delta}{2^{N - 1}} < \veps \quad \text{for all } \bm{\gamma} \in I_{N},
    \]
    which implies that $I_{N} \subseteq B_{\veps}(\bm{\beta}) \subseteq G_{\alpha_{0}}$. Therefore, $I_{N}$ has a
    finite subcover $\set*{G_{\alpha_{0}}}$, which is a contradiction. Thus $I$ has a finite subcover, and $I$ is compact.
\end{proof}

\bigskip

\begin{corollary} \label{cor:ClosedIntervalInRIsCompact}
    For any $- \infty < a \leq b < + \infty$, the closed interval $[a, b]$ is compact in $\paren*{\R, \norm*{\cdot}}$.
\end{corollary}

\bigskip

\begin{theorem} \label{thm:HeineBorelTheorem}
    Let $\paren*{\R^{n}, \norm*{\cdot}}$ be given for some $n \in \N$ and $E \subseteq \R^{n}$.
    Then we have
    \[
    E \text{ is compact} \iff E \text{ is closed and bounded}.
    \]
    This theorem is called the \textbf{Heine-Borel Theorem}.
\end{theorem}

\begin{proof}
    \textbf{($\implies$)} Let $E$ be compact. $E$ is closed by \Cref{thm:CompactSetIsClosed}, and $E$ is
    bounded by \Cref{thm:CompactSetIsBounded}.
    \\[6mm]
    \textbf{($\impliedby$)} Let $E$ be closed and bounded. Then there exists some $\bm{p} \in \R^{n}$ and
    $r \in \R^{+}$ such that $E \subseteq B_{r}(\bm{p})$. Let $I$ be a $n$-cell in $\R^{n}$ defined by
    \[
    I \coloneqq \prod_{j = 1}^{n} \bracket*{p_{j} - r, p_{j} + r} \subseteq \R^{n}.
    \]
    Then we have $E \subseteq B_{r}(\bm{p}) \subseteq I$. By \Cref{thm:KCellIsCompact}, $I$ is compact.
    Since $E$ is closed, by \Cref{thm:ClosedSubsetOfCompactSetIsCompact}, $E$ is compact.
\end{proof}

\end{section}